\section{LoRA-Based Behavioural Adaptation}
\label{sec:lora}

The preceding modules manage \emph{what} legal knowledge the system uses; this section describes \emph{how} the model reasons over that knowledge. To balance the otherwise external-memory-heavy design, behavioural competence is acquired through parameter-efficient fine-tuning while the base model remains frozen, in line with the hybrid framing established in the literature review.

\subsection{Role and the Knowledge-Behaviour Separation}
A guiding principle separates two concerns: \emph{retrieval decides which statute to use, whereas the adapter decides how to reason with it}. Low-Rank Adaptation (LoRA) \parencite{hu2021loralowrankadaptationlarge} is used to learn behaviour and policy---reading the retrieved context, producing legally structured reasoning (e.g. the major-premise / minor-premise / conclusion form of a syllogism), citing the relevant articles, and respecting the output constraints of the unlearning module---rather than to memorize legal facts, which remain in the external knowledge base.

\subsection{Why Adapters for Behaviour, Not Facts}
This separation is justified by the singular-value analysis of \textcite{biderman2024loralearnsforgets}: the intrinsic rank of the weight change required for continual \emph{pre-training} is extremely high (on the order of $2000$), so low-rank adapters cannot absorb large volumes of factual knowledge; in contrast, \emph{instruction} fine-tuning is well approximated by a low intrinsic rank, where LoRA matches full fine-tuning. Storing facts in adapters would therefore be both under-capacity and forgetting-prone, whereas LoRA used for behaviour acts as a strong regularizer that ``learns less and forgets less'', preserving the base model's general capabilities.

\subsection{Adapter Configuration and Training Data}
The proposed configuration is summarized in Table~\ref{tab:lora_config}. The base model is an instruction-tuned decoder-only LLM, chosen because such architectures resist forgetting and provide a protective instruction-tuning buffer \parencite{li2024revisitingcatastrophicforgettinglarge}. Consistent with the low-resource setting that motivates this work, all experiments use \emph{small} open-weight models---primarily in the $1$--$4$B-parameter range and at most ${\sim}15$B---rather than frontier-scale systems. This is a deliberate design choice as much as a hardware constraint: because the base remains frozen and all factual knowledge lives in the external store, the framework is intended to deliver compliant, updatable legal assistance \emph{without} relying on large parametric capacity, and a small frozen base is precisely the regime where the cost and forgetting advantages of the RAG-centric design matter most. The same small scale is used for the parametric comparison baselines so that the forgetting analysis is a like-for-like comparison rather than a confound of model size. To preserve general ability, a small fraction (about $0.5\%$) of general Vietnamese text is mixed into the instruction data as rehearsal \parencite{scialom-etal-2022-fine}. The instruction set comprises legal QA, NLI, and syllogism examples formatted with task-distinct prompt templates to avoid label-space collision \parencite{yin-etal-2022-contintin}, and includes ``read-the-retrieved-context-then-answer'' examples so that the adapter explicitly learns the policy of grounding on the legal RAG memory.

\begin{table}[H]
\centering
\small
\caption{Proposed LoRA configuration for behavioural adaptation.}
\label{tab:lora_config}
\begin{tabularx}{\textwidth}{l >{\raggedright\arraybackslash}X >{\raggedright\arraybackslash}X}
\hline
\textbf{Component} & \textbf{Setting} & \textbf{Rationale} \\
\hline
Base model & Instruction-tuned decoder-only, frozen ($1$--$4$B, ${\le}15$B) & forgetting resistance + SFT buffer; fits low-resource budget \\
Rank $r$ & 256 & headroom for multi-task behavioural coverage \\
Scaling $\alpha$ & 512 ($\alpha = 2r$) & standard scaling \\
Target modules & all linear projections & full behavioural coverage \\
Rehearsal & $\approx 0.5\%$ general Vietnamese text & preserves zero-shot ability \\
Optimizer & AdamW (optionally SAM) & flat minima reduce forgetting \\
\hline
\end{tabularx}
\end{table}

\subsection{Sequential Multi-Sector Adaptation}
If adapters are trained sequentially across legal sub-domains---in this study the sequential test scenario adapts on \textbf{labour law} and then \textbf{tax law}, two legislation-rich sub-domains chosen to make cross-domain forgetting observable---naive sharing of a single adapter causes parameter collision, and adapting to the later sub-domain can degrade performance on the earlier one (catastrophic forgetting, measured here through Backward Transfer). Two remedies from the literature are applicable: orthogonal-subspace projection, as in O-LoRA \parencite{wang-etal-2023-orthogonal} and its sparsity-based refinement N-LoRA \parencite{yang2024parametercollisionhinderingcontinual}, which isolate task updates; and soft expert routing, as in SLIM \parencite{han-etal-2025-slim}, which routes queries of different sub-domains to dedicated experts. Orthogonal methods are noted to be sensitive to their projection hyperparameters and are tuned accordingly. This two-domain sequence is the minimal configuration that makes cross-domain forgetting measurable; longer sequences over further sub-domains are a straightforward extension.

\subsection{Interaction with Retrieval at Inference}
At inference, the frozen base model augmented by the LoRA adapter receives the selectively retrieved and unlearning-filtered context, the user query, and the output constraint $Q$, and produces a grounded legal argument, as shown in Figure~\ref{fig:lora_inference}.

\begin{figure}[H]
\centering
\begin{tikzpicture}[
  font=\footnotesize, node distance=0.7cm and 1.2cm,
  b/.style={rectangle, rounded corners, draw, align=center, text width=3.0cm, minimum height=0.8cm},
  frz/.style={rectangle, rounded corners, draw, align=center, text width=2.6cm, fill=blue!6},
  lora/.style={rectangle, rounded corners, draw, very thick, align=center, text width=2.6cm, fill=green!10},
  arr/.style={-{Stealth[length=2mm]}, thick}
]
\node[b] (ctx) {Retrieved context\\ (selective + unlearned)};
\node[b, below=of ctx] (q) {Query $q$ + constraint $Q$};
\node[frz, right=of ctx] (base) {Frozen base LLM};
\node[lora, below=of base] (lo) {LoRA adapter\\ (behaviour, IFT)};
\node[b, right=of base] (out) {Legal reasoning\\ (syllogism + citation)};
\draw[arr] (ctx) -- (base);
\draw[arr] (q) -- (lo);
\draw[arr] (base) -- (out);
\draw[arr] (lo) -- (out);
\end{tikzpicture}
\caption{Inference-time interaction between the frozen base model, the behavioural LoRA adapter, and the retrieved legal context.}
\label{fig:lora_inference}
\end{figure}

In summary, the LoRA layer supplies the parametric, behavioural counterpart to the non-parametric memory modules: it teaches the model to reason and cite in the legal register without ever storing volatile facts in its weights, keeping the system's knowledge updatable and unlearnable through retrieval alone.
